\chapter{The Training Procedure}
\chaptermark{Training}
\label{ch:training}

\section{The full loss of one batch}

A batch has 32 samples (about 5\% of them non-speech). With the label conventions of Chapters~\ref{ch:whisper} and
\ref{ch:robust}:
\[
\mathcal L=\underbrace{\frac{\sum_{i,t}\omega_{it}\,\1[y_{it}\ne-100]\,\big(-\log p_\theta(y_{it}\mid\cdot)\big)}{\sum_{i,t}\1[y_{it}\ne-100]}}_{\text{token cross-entropy}}
\;+\;\lambda\underbrace{\frac{1}{|\mathcal S|}\sum_{i\in\mathcal S}\frac{\ell_{\text{CTC},i}}{T_i}}_{\text{LID-CTC}},
\]
with $\omega_{it}=1$ except $\omega_{i,3}=w_{\text{ns}}$ on the EOT of a non-speech sample ($w_{\text{ns}}=1$ in Stage 2), and
$\lambda=0.3$.

\section{The optimiser}

\begin{definition}{AdamW}{adamw}
For each parameter $\theta$ with gradient $g_s$ at step $s$:
\begin{align*}
m_s&=\beta_1m_{s-1}+(1-\beta_1)g_s, & v_s&=\beta_2v_{s-1}+(1-\beta_2)g_s^2,\\
\hat m_s&=m_s/(1-\beta_1^s), & \hat v_s&=v_s/(1-\beta_2^s),\\
\theta&\leftarrow\theta-\eta_s\Big(\frac{\hat m_s}{\sqrt{\hat v_s}+\epsilon}+\lambda_{\text{wd}}\theta\Big).
\end{align*}
Here $\beta_1=0.9$, $\beta_2=0.98$, $\epsilon=10^{-6}$; weight decay $0$ for LoRA, the embedding and the head.
\end{definition}

The bias corrections $1/(1-\beta^s)$ undo the bias toward $0$ caused by initialising the moments at $0$. Dividing by $\sqrt{\hat v}$
gives every parameter a step of roughly size $\eta$ regardless of its gradient scale.

Three parameter groups with different learning rates:
\begin{center}
\begin{tabular}{lcc}
\toprule
Group & peak learning rate & weight decay\\
\midrule
LoRA matrices (22 encoder + 4 decoder layers) & $3\times10^{-4}$ & 0\\
Embedding (only the \tok{or} row can change) & $2.5\times10^{-4}$ & 0\\
LID-CTC head & $10^{-3}$ & 0\\
\bottomrule
\end{tabular}
\end{center}

\section{Learning-rate schedule}

\begin{definition}{Linear warm-up, linear decay}{sched}
With $S$ total steps and $W=\max(1,0.1S)$ warm-up steps, every group's learning rate is its peak times
\[
\phi(s)=\begin{cases}\dfrac{s+1}{W}, & s<W,\\[6pt] \max\!\Big(0,\ \dfrac{S-s}{\max(1,S-W)}\Big), & s\ge W.\end{cases}
\]
\end{definition}
For Stage 2, $S=11{,}000$ and $W=1{,}100$. Warm-up matters for Adam: in the first steps $\hat v$ is estimated from very few gradients
and can be tiny, which would make the first updates huge.

\section{Gradient clipping and precision}

\paragraph{Clipping.} If the global gradient norm $\|g\|=\big(\sum_{\text{all params}}\|g_\theta\|^2\big)^{1/2}$ exceeds 1, all
gradients are rescaled by $1/\|g\|$. This bounds the size of any single update (a single bad batch cannot throw the model far away).

\paragraph{bfloat16 autocast.} Matrix products run in bfloat16 (8-bit exponent like float32, 7-bit mantissa), halving memory traffic
with the same numerical range as float32; the loss and softmax are computed in float32.

\paragraph{\texttt{torch.compile}.} The forward pass is compiled into fused kernels; it gave 12.7 $\to$ about 23 utterances/s in a
benchmark (15--17 in real Whisper-small runs). To limit the number of distinct shapes that must be compiled, decoder sequences are
padded to a multiple of 64 tokens.

\section{Reproducible randomness and resumable training}

\begin{definition}{Seeded data order}{seeded}
Epoch $e$ uses the permutation \texttt{RandomState(seed}$+e$\texttt{).permutation(n)}; batch $b$ is positions $32b..32b+31$. Sample
$i$ in epoch $e$ draws all its augmentation randomness from \texttt{Random(f"\{seed\}:\{e\}:\{i\}")}. SpecAugment has its own
generator.
\end{definition}

Because every random choice is a function of (seed, epoch, index), the data a run sees does not depend on how many DataLoader
workers there are or in which order they finish, and a resumed run sees exactly the same samples as an uninterrupted one.

\begin{incode}
A checkpoint (every 500 steps, written atomically via a temporary file and \texttt{os.replace}) stores: trainable parameters,
head, optimiser and scheduler states, step, epoch and batch position, the SpecAugment and CUDA random states (CUDA RNG drives LoRA
dropout), loss history, partial results and elapsed time. \texttt{--resume} restores all of it. Verified: resuming 20 $\to$ 40
steps differs from an uninterrupted run by 0.004 in loss, less than two identical fresh runs differ from each other (0.009--0.014,
GPU non-determinism).
\end{incode}

\section{Stage 2 in numbers}

\begin{center}
\begin{tabular}{ll}
\toprule
Base model & Whisper large-v3-turbo\\
Data & 700 h mix (349{,}592 references), 7 languages + MUCS\\
Steps $\times$ batch & $11{,}000\times32=352{,}000$ samples $\approx$ 1 epoch\\
LoRA & $r=32$, $\alpha=64$, dropout 0.05, encoder layers 0--9 frozen\\
Extras & \tok{or} row, LID-CTC $\lambda=0.3$, augmentation, 5\% non-speech\\
Throughput & 5.9 utterances/s; 59{,}895 s $=$ 16.6 h; peak memory 88.6 GB\\
Final training loss & 0.159 (mean of last 50 steps)\\
\bottomrule
\end{tabular}
\end{center}

\begin{figure}[h]
\centering
\begin{tikzpicture}
\begin{axis}[width=0.9\linewidth,height=6cm,xlabel={training step},ylabel={loss (20-step mean)},
  xmin=0,xmax=11000,ymin=0,ymax=0.5,grid=both,grid style={gray!20},legend pos=north east,
  scaled x ticks=false,xticklabel style={/pgf/number format/.cd,1000 sep={\,}}]
\addplot[thick,color=titleblue,mark=*,mark size=1.2pt] coordinates {
(500,0.4250)(1000,0.3580)(1500,0.3193)(2000,0.2939)(2500,0.2776)(3000,0.2555)(3500,0.2415)(4000,0.2487)(4500,0.2259)
(5000,0.2235)(5500,0.2152)(6000,0.2073)(6500,0.2001)(7000,0.1985)(7500,0.1919)(8000,0.1895)(8500,0.1779)(9000,0.1733)
(9500,0.1669)(10000,0.1723)(10500,0.1771)(11000,0.1622)};
\addlegendentry{total loss}
\draw[dashed,gray] (axis cs:1100,0) -- (axis cs:1100,0.5) node[pos=0.92,right,font=\footnotesize]{end of warm-up};
\end{axis}
\end{tikzpicture}
\caption{Stage 2 training loss (from \file{logs/s2_turbo.log}; step 1 loss was 2.13, off the chart). The LID-CTC term fell from 2.04
at step 1 to about 0.003 by step 1{,}000 and 0.001 at the end: the head learns quickly, and afterwards almost all of the loss is
transcription.}
\end{figure}

\begin{figure}[h]
\centering
\begin{tikzpicture}
\begin{axis}[width=0.9\linewidth,height=5.5cm,xlabel={training step},ylabel={mean dev CER (19 sets)},
  xmin=0,xmax=11500,ymin=0.08,ymax=0.21,grid=both,grid style={gray!20},axis y line*=left,
  scaled x ticks=false,xticklabel style={/pgf/number format/.cd,1000 sep={\,}}]
\addplot[thick,color=thmc,mark=square*,mark size=1.4pt] coordinates {
(1000,0.1997)(2000,0.1495)(3000,0.1328)(4000,0.1371)(5000,0.1234)(6000,0.1218)(7000,0.1102)(8000,0.1077)(9000,0.1045)(10000,0.1068)(11000,0.1004)};
\label{plot:cer}
\end{axis}
\begin{axis}[width=0.9\linewidth,height=5.5cm,xmin=0,xmax=11500,ymin=0.75,ymax=0.95,axis y line*=right,axis x line=none,
  ylabel={mean dev token-LID},legend pos=north east,legend style={font=\footnotesize}]
\addlegendimage{/pgfplots/refstyle=plot:cer}\addlegendentry{dev CER (left)}
\addplot[thick,color=exc,mark=triangle*,mark size=1.6pt,dashed] coordinates {
(1000,0.771)(2000,0.861)(3000,0.906)(4000,0.901)(5000,0.911)(6000,0.932)(7000,0.923)(8000,0.912)(9000,0.929)(10000,0.930)(11000,0.931)};
\addlegendentry{dev token-LID (right)}
\end{axis}
\end{tikzpicture}
\caption{Dev-set evaluation inside the Stage 2 run (48 utterances per set, \file{runs/s2_turbo/result.json}). CER was still falling at
the last evaluation (0.105 at step 9{,}000, 0.100 at the end), while LID plateaued around 0.93 after 6{,}000 steps.}
\end{figure}

\begin{exercise}{Learning rate at step 6050}{lrcalc}
\textit{Solution.} $s=6050\ge W=1100$: $\phi=(11000-6050)/(11000-1100)=4950/9900=0.5$. LoRA learning rate $1.5\times10^{-4}$, Odia row
$1.25\times10^{-4}$, head $5\times10^{-4}$.
\end{exercise}

\begin{vivatip}
Why only one epoch? Compute: one epoch is 16.6 h on the only GPU, and the plan was to evaluate before spending more. Note also that
the dev CER evaluated every 1{,}000 steps was still improving slowly at the end, so more training (on a reseeded mix) might help;
this is listed as future work (see the dev curve above: 0.200 at step 1{,}000, 0.100 at the end, not yet flat).
\end{vivatip}
